\subsection{\texorpdfstring{特征 \(p \neq 0\) 下的可积丛}{特征 p 非零时的可积丛}}

本节中，假设 K 的特征为 \(p \neq 0\)。记 X 为一个局部有限维的 K-解析流形。

9.4.1（“p 次幂”）。设 \(\xi\) 是 X 的一个开子集上的向量场。在 U 上存在且仅存在一个向量场 \(\xi^p\)，使得

\[
\mathrm{D} _ {\xi^ {p}} (f) = (\mathrm{D} _ {\xi}) ^ {p} (f) = \underbrace {\mathrm{D} _ {\xi} (\mathrm{D} _ {\xi} (\dots (\mathrm{D} _ {\xi} (f)) . . .))} _ {p \text { times }}\tag{1}
\]

对定义在 U 的一个开子集上的每个解析函数 f，均有

若 \(\varphi\) 是 U 上的解析函数，则有

\[
(\varphi \xi) ^ {p} = \varphi^ {p} \xi^ {p} + (\mathrm{D} _ {\varphi \xi}) ^ {p - 1} (\varphi). \xi\tag{2}
\]

若 \(\eta\) 是 U 上的向量场，则有

\[
[ \xi^ {p}, \eta ] = \operatorname{ad}(\xi) ^ {p} (\eta) = [ \xi , [ \xi , \dots , [ \xi , \eta ] \dots ] ].\tag{3}
\]

9.4.2（“Jacobson 恒等式”）。设 L 是自由 \(F_{p}\)-李代数（LIE, II, § 2, no. 2），它以含有两个元素的集合 \(\{x, y\}\) 为生成集合。L 中存在且仅存在一个元素 \(\Lambda_{p}(x, y)\)，使得

\[
(x + y) ^ {p} = x ^ {p} + y ^ {p} + \Lambda_ {p} (x, y)\tag{4}
\]

该等式在 L 的包络代数中成立。例：

\[
\Lambda_ {2} (x, y) = [ x, y ]; \quad \Lambda_ {3} (x, y) = [ x, [ x, y ] ] - [ y, [ x, y ] ].
\]

若 \(\xi\) 和 \(\eta\) 是 X 上的向量场，则有

\[
(\xi + \eta) ^ {p} = \xi^ {p} + \eta^ {p} + \Lambda_ {p} (\xi , \eta).\tag{5}
\]

9.4.3. 例。--- 取 \(X = K\)（分别为 \(K^*\)），并记 x 为典范映射 \(X \to K\)。设 \(\xi\)（分别为 \(\eta\)）是向量场 \(\partial/\partial x\)（分别为 \(x\cdot\partial/\partial x\)）；它在 X 的加法群（分别为乘法群）的平移下不变。我们有

\[
\xi^ {p} = 0 \quad \text { and } \quad \eta^ {p} = \eta .\tag{6}
\]

9.4.4（“可积丛”）。设 F 是 T(X) 的一个向量子丛。称 F 可积，如果对于每个 \(x \in X\)，存在 X 在 x 处的一个坐标系（\(\zeta^{1}, \ldots, \zeta^{n}\)）以及一个整数 \(m \leqslant n\)，使得 \((\partial/\partial\zeta^{1}, \ldots, \partial/\partial\zeta^{m})\) 是 F 在 x 的某个邻域内的标架。

F 可积的充要条件是满足以下条件：

对于 X 的每个开集 U，\(\mathcal{S}_{\mathbb{F}}^{\omega}(\mathbf{U})\) 是 \(\mathcal{S}_{\mathrm{T}(\mathbf{X})}^{\omega}(\mathbf{U})\) 的李子代数，并在运算 \(\xi \mapsto \xi^{p}\) 下稳定。

